\section{LLMs and Prompting}
\label{sec:llms}

Large Language Models are typically pretrained to predict tokens conditioned on
preceding context. This objective supports general-purpose instruction following
after further training, enabling tasks such as summarisation, translation,
question answering, and code generation
\parencites{brown2020language}{chen2025struq}.

In an LLM-integrated application, the model may receive application-provided
higher-priority instructions, a user request, retrieved data, and tool results.
Interfaces normally encode these components using role markers or delimiters,
and instruction-hierarchy training can teach a model to prioritise some channels
over others
\parencite{wallace2024instruction}. These distinctions are meaningful inputs to
the model. They do not, however, constitute an independently enforced authority
mechanism: generation still depends on learned behaviour over the combined
context.

Prompt injection exploits failures of that learned prioritisation. An attacker
crafts text in an untrusted channel so that the model treats it as a controlling
instruction rather than as task-relevant data
\parencite{liu2023promptinjection}. Structured prompting can make the intended
roles clearer and model training can improve adherence to them, but downstream
security should not assume that formatting alone prevents untrusted text from
influencing generation.
